\subsection{OCTONIONS}

Let E be a quaternion algebra of type \((\alpha, \beta, \gamma)\) over A (§2, no. 5, Example 2) and let \(\delta \in A\) . The Cayley extension F of E by \(\delta\) and the conjugation of E is called an octonion algebra over A and is said to be of type \((\alpha, \beta, \gamma, \delta)\) . By Proposition 3 of no. 2, F is an alternative algebra. It has a basis \((e_{i})_{0 \leq i \leq 7}\) of 8 elements, defined by

\[
\begin{array}{l l l l} e _ {0} = (e, 0), & e _ {1} = (i, 0), & e _ {2} = (j, 0), & e _ {3} = (k, 0) \\ e _ {4} = (0, e), & e _ {5} = (0, i), & e _ {6} = (0, j), & e _ {7} = (0, k) \end{array}
\]

where \((e, i, j, k)\) is the basis of E defined loc. cit.; clearly \(e_0\) (also denoted by \(e\) ) is the unit element of F. If \(u = \sum_{i=0}^{7} \xi_i e_i\) is an element of F (with the \(\xi_i \in A\) ), formulae (23), (24) and (31) of § 2, no. 5, give for the conjugate, trace and norm of the octonion \(u\)

\[
\left\{ \begin{array}{c} \bar {u} = (\xi_ {0} + \beta \xi_ {1}) e _ {0} - \sum_ {i = 1} ^ {7} \xi_ {i} e _ {i} \\ \mathrm{T} _ {\mathrm{F}} (u) = 2 \xi_ {0} + \beta \xi_ {1} \\ \mathrm{N} _ {\mathrm{F}} (u) = \xi_ {0} ^ {2} + \beta \xi_ {0} \xi_ {1} - \alpha \xi_ {1} ^ {2} - \gamma (\xi_ {2} ^ {2} + \beta \xi_ {2} \xi_ {3} - \alpha \xi_ {3} ^ {2}) \\ - \delta (\xi_ {4} ^ {2} + \beta \xi_ {4} \xi_ {5} - \alpha \xi_ {5} ^ {2}) + \gamma \delta (\xi_ {6} ^ {2} + \beta \xi_ {6} \xi_ {7} - \alpha \xi_ {7} ^ {2}). \end{array} \right.\tag{4}
\]

Now let \(u = (x, y)\) , \(u' = (x', y')\) and \(u'' = (x'', y'')\) be three octonions (where the elements \(x, x', x'', y, y', y''\) belong to E). Formulae (24) and (27) of § 2, no. 5 give

\[
\begin{array}{r l} & {\mathrm{T} _ {\mathrm{F}} ((u u ^ {\prime}) u ^ {\prime \prime}) = \mathrm{T} (x x ^ {\prime} x ^ {\prime \prime}) + \delta \mathrm{T} (\bar {y} ^ {\prime} y x ^ {\prime \prime}) + \delta \mathrm{T} (\bar {y} ^ {\prime \prime} y \bar {x} ^ {\prime}) + \delta \mathrm{T} (\bar {y} ^ {\prime \prime} y ^ {\prime} x)} \\ & {\mathrm{T} _ {\mathrm{F}} (u (u ^ {\prime} u ^ {\prime \prime})) = \mathrm{T} (x x ^ {\prime} x ^ {\prime \prime}) + \delta \mathrm{T} (x ^ {\prime \prime} \bar {y} ^ {\prime} y) + \delta \mathrm{T} (\bar {x} ^ {\prime} \bar {y} ^ {\prime \prime} y) + \delta \mathrm{T} (x \bar {y} ^ {\prime \prime} y ^ {\prime})} \end{array}
\]

(where T denotes trace in E and use is made of the fact that E is associative). As \(\mathrm{T}(xy) = \mathrm{T}(yx)\) for all quaternions x, y (§ 2, no. 4, formula (17)), it follows that

\[
\mathrm{T} _ {\mathrm{F}} ((u u ^ {\prime}) u ^ {\prime \prime}) = \mathrm{T} _ {\mathrm{F}} (u (u ^ {\prime} u ^ {\prime \prime})).\tag{5}
\]

We study in particular octonions of type \((-1, 0, -1, -1)\) ; formulae (4) then simplify to

\[
\left\{ \begin{array}{c} {\bar {u}} = \xi_ {0} e _ {0} - \sum_ {i = 1} ^ {7} \xi_ {i} e _ {i} \\ \mathrm{T} _ {\mathrm{F}} (u) = 2 \xi_ {0} \\ \mathrm{N} _ {\mathrm{F}} (u) = \sum_ {i = 0} ^ {7} \xi_ {i} ^ {2}. \end{array} \right.\tag{6}
\]

*If we take A to be the field R of real numbers, the octonions of type \((-1,0,-1,-1)\) over R are called Cayley octonions (or octaves). It follows from Proposition 2 (ii) of no. 2 that every Cayley octonion \(\neq0\) is invertible.*

PROPOSITION 4. Let F be an octonion algebra of type \((-1,0,-1,-1)\) over A. There exists a vector space V of dimension 3 over the field with two elements Z/2Z and a bijection \(\lambda\mapsto e_{\lambda}^{\prime}\) of V onto the basis \((e_{i})_{0\leq i\leq7}\) such that

\[
e _ {0} ^ {\prime} = e _ {0}, \quad e _ {\lambda} ^ {\prime} e _ {\mu} ^ {\prime} = \pm e _ {\lambda + \mu} ^ {\prime}\tag{7}
\]

for all \(\lambda, \mu\) in V. In order that \(e_{\lambda}'(e_{\mu}'e_{\nu}') = (e_{\lambda}'e_{\mu}')e_{\nu}'\) , it suffices that, in V, \(\lambda, \mu, \nu\) be linearly dependent over Z/2Z; this condition is necessary if \(2 \neq 0\) in A.

We preserve the notation at the beginning of this no. It follows from formulae (33) of §2, no. 5 that the set S consisting of the elements \(\pm e_0\) , \(\pm e_1\) , \(\pm e_2\) , \(\pm e_3\) is stable under multiplication. Moreover, for \(x, y, y'\) in E, by formula (22) of §2, no. 5,

\[
(x, 0) (0, y ^ {\prime}) = (0, y ^ {\prime} x), \quad (0, y ^ {\prime}) (x, 0) = (0, y ^ {\prime} \bar {x}), \quad (0, y) (0, y ^ {\prime}) = (- \bar {y} ^ {\prime} y, 0)\tag{8}
\]

so that the set T consisting of the elements \(\pm e_{i}\) ( \(0 \leqslant i \leqslant 7\) ) is stable under multiplication; moreover, its multiplication table is independent of the ring A.

In particular, let \(\mathbf{A}''\) be the field \(\mathbf{Z}/2\mathbf{Z}\) with two elements and let \(\mathbf{E}''\) be the quaternion algebra of type (1, 0, 1) over \(\mathbf{A}''\) and \(\mathbf{F}''\) the algebra of octonions of type (1, 0, 1, 1) over \(\mathbf{A}''\) ; let \((e_i'')_{0 \leq i \leq 7}\) be the basis of \(\mathbf{F}''\) formed as described above. Since \(-e_i'' = e_i''\) , the set \(\mathbf{T}''\) of \(e_i''\) has 8 elements and is stable under multiplication; moreover, it follows immediately from the above that the mapping \(\theta: \mathbf{T} \to \mathbf{T}''\) such that \(\theta(e_i) = \theta(-e_i) = e_i''\) for \(0 \leq i \leq 7\) is a homomorphism for multiplication. Moreover the quaternion algebra \(\mathbf{E}''\) is in this case commutative and hence \(\mathbf{F}''\) is associative ( \(\S 2\) , no. 5, Proposition 5); further, conjugation in \(\mathbf{F}''\) is in this case the identity. Hence \(\mathbf{T}''\) is a group and formulae (8) show that it is commutative; these formulae and formulae (33) of \(\S 2\) , no. 5 show that the square of every element of \(\mathbf{T}''\) is the unit. If V is used to denote the group \(\mathbf{T}''\) written additively, V can be given a unique vector space structure over \(\mathbf{Z}/2\mathbf{Z}\) , necessarily of dimension 3 since

\[
\operatorname{Card} (\mathbf {V}) = 8 = 2 ^ {\dim (\mathbf {V})}.
\]

For all \(\lambda \in V\) , let \(e_{\lambda}'\) then denote the element of \((e_i)_{0 \leq i \leq 7}\) such that \(\theta(e_{\lambda}') = \dot{\lambda}\) ; then \(e_0' = e_0\) ; moreover, as \(\theta\) is a homomorphism and the relation \(\theta(x) = \theta(y)\) is equivalent to \(x = \pm y\) , \(e_{\lambda}' e_{\mu}' = \pm e_{\lambda + \mu}'\) . If \(\lambda, \mu, \nu\) are linearly independent over \(\mathbf{Z}/2\mathbf{Z}\) , they form a basis of \(V\) and hence all the elements \(e_i\) ( \(0 \leqslant i \leqslant 7\) ) would belong to the subalgebra generated by \(e_{\lambda}', e_{\mu}'\) and \(e_{\nu}'\) ; when \(2 \neq 0\) in \(A\) , \(e_{\lambda}'(e_{\mu}' e_{\nu}') = (e_{\lambda}' e_{\mu}') e_{\nu}'\) is therefore impossible, for \(F\) would be associative and hence \(E\) would be commutative ( \(\S 2\) , no. 5, Proposition 5), which contradicts relations (33) of \(\S 2\) , no. 5. On the other hand, if \(\lambda, \mu, \nu\) are linearly

dependent in V, the three elements \(e_{\lambda}^{\prime}\) , \(e_{\mu}^{\prime}\) , \(e_{\nu}^{\prime}\) belong to a subalgebra with 2 generators of F, which is therefore associative (no. 1, Proposition 1); whence the conclusion.

Remark. As \(\bar{e}_{\lambda}^{\prime} = -e_{\lambda}^{\prime}\) for \(\lambda \neq 0\) ,

\[
e _ {\lambda} ^ {\prime 2} = - e _ {0} \quad \text { for } \lambda \neq 0,
\]

\[
e _ {\mu} ^ {\prime} e _ {\lambda} ^ {\prime} = - e _ {\lambda} ^ {\prime} e _ {\mu} ^ {\prime} \quad \text { for } \lambda \neq 0, \mu \neq 0 \text { and } \mu \neq \lambda .
\]

